\documentclass[11pt,oneside]{article}
\input{../../../Header/packages.sty}
\input{../../../Header/mathexpressions.sty}
\usepackage{fancyhdr}
\pagestyle{fancy}
\fancyhf{}
\lhead{Lektionsplanung: Der Höhensatz}
\rhead{Jonas Guler (Gu)}
\rfoot{\thepage}
\renewcommand{\headrulewidth}{0.4pt}
\setlength{\headheight}{14pt}
\setlength{\parindent}{0pt}
\setlength{\parskip}{4pt}
\linespread{1.15}

\begin{document}
{\Large\textbf{Lektionsplanung: Der Höhensatz}}\\[2mm]
\begin{tabularx}{\textwidth}{@{}lX}
	Thema: & Sätze im rechtwinkligen Dreieck, Höhensatz\\
	Dauer: & 1 Lektion à 45 Minuten\\
	Einbettung: & nach dem Satz des Pythagoras und dem Kathetensatz (Skript \textit{Sätze im rechtwinkligen Dreieck}, Abschnitt 3, und Lektion \textit{Der Kathetensatz})\\
	Material: & Arbeitsblatt \texttt{höhensatz.pdf}, Taschenrechner, Geodreieck und Zirkel\\
\end{tabularx}

\section*{Lernziele}
Die Lernenden \ldots
\begin{itemize}[itemsep=0pt]
	\item entdecken durch Konstruktion und an Zahlenbeispielen den Zusammenhang $h_c^2=p\cdot q$ und formulieren ihn als Vermutung.
	\item können den Höhensatz in Worten und als Formel angeben und am Flächenbild (Höhenquadrat = Rechteck aus $p$ und $q$) erklären.
	\item können den Höhensatz aus Pythagoras und Kathetensatz herleiten und kennen die ganze Satzgruppe des Pythagoras.
	\item berechnen mit dem Höhensatz und den bekannten Sätzen fehlende Grössen im rechtwinkligen Dreieck und wenden den Satz bei Flächenverwandlungen an.
\end{itemize}

\textbf{Vorwissen:} Satz des Pythagoras, Kathetensatz $a^2=c\cdot q$, $b^2=c\cdot p$, Satz des Thales, Konstruktion einer Senkrechten.

\textbf{Notation:} Wie im Skript und in der Kathetensatz-Lektion liegt $p$ bei $A$ (an der Kathete $b$) und $q$ bei $B$. Der Höhensatz selbst ist symmetrisch in $p$ und $q$, die Notation wird erst bei Rechnungen mit dem Kathetensatz wichtig.

\section*{Ablauf}
\renewcommand{\arraystretch}{1.3}
\begin{tabularx}{\textwidth}{|>{\raggedright\arraybackslash}p{11mm}|>{\raggedright\arraybackslash}p{27mm}|X|>{\raggedright\arraybackslash}p{22mm}|}\hline
	\textbf{Zeit} & \textbf{Phase} & \textbf{Inhalt und Aktivität} & \textbf{Sozialform}\\\hline
	0'--3' & Einstieg & Begrüssung, kurze Repetition Kathetensatz (Merksatz abfragen). Lektionsziel: \textit{Nach den Katheten untersuchen wir heute die Höhe.} Arbeitsblatt austeilen. & Plenum\\\hline
	3'--15' & Schüleraktivität (Entdeckung) & \textbf{Aufgabe 1.1, Teil A:} Dreieck aus $p=4.5$, $q=2$ mit Thaleskreis konstruieren, $h_c$ messen, Tabelle ausfüllen, Vermutung formulieren. \textbf{Teil B} (Beweis nur mit Pythagoras) für schnelle Paare. Häufige Hürde: Wie findet man $C$? Hinweis auf Thaleskreis über $\overline{AB}$ und Senkrechte in $H_c$. & Partnerarbeit\\\hline
	15'--26' & Theorieinput & Vermutungen sammeln. \textbf{Satz 2.1} mit Flächenbild an der Wandtafel (Quadrat über $h_c$, Rechteck $p\times q$, $q$ heruntergeklappt). \textbf{Beweis} gemeinsam mit Pythagoras im Teildreieck und Kathetensatz; Teil B als zweiten Beweis von einem Paar vorstellen lassen. \textbf{Bemerkungen:} geometrisches Mittel, Symmetrie in $p$ und $q$, Übersichtstabelle Satzgruppe gemeinsam ausfüllen. \textbf{Beispiel 1} gemeinsam lösen (ca.\ 3'). & Plenum, Lehrervortrag mit Fragen\\\hline
	26'--44' & Übungen & \textbf{Pflicht:} Aufgabe 3.1 a)--d), 3.2, 3.3 a)--d).\newline \textbf{Zusatz:} 3.1 e)--f), 3.3 e), 3.4 (Wurzeln konstruieren), 3.5 (Tunnel), 3.6 $\bigstar$ (Mittelwerte). Lösungen liegen bei der Lehrperson auf bzw.\ werden ab Min.\ 38 aufgelegt. & Einzel- oder Partnerarbeit\\\hline
	44'--45' & Abschluss & Merksatz wiederholen lassen: \textit{Höhe$^2$ = Produkt der Hypotenusenabschnitte.} Hausaufgabe: Pflichtaufgaben fertig lösen. & Plenum\\\hline
\end{tabularx}

\section*{Tafelbild (Vorschlag)}
\begin{minipage}{0.52\linewidth}
	\textbf{Höhensatz}\\
	Höhenquadrat = Rechteck aus den beiden Hypotenusenabschnitten:
	\[h_c^2=p\cdot q\qquad\Leftrightarrow\qquad h_c=\sqrt{p\cdot q}\]
	\textbf{Beweis:} (I) $h_c^2=b^2-p^2$ (Pythagoras in $\triangle AH_cC$),\\ (II) $b^2=c\cdot p$ (Kathetensatz).\\
	$h_c^2=c\cdot p-p^2=p\,(c-p)=p\cdot q$.
\end{minipage}
\hfill
\begin{minipage}{0.42\linewidth}
	\centering
	\begin{tikzpicture}[x=1cm, y=1cm, font=\footnotesize, scale=0.6]
		\filldraw[fill=myblue, fill opacity=0.25, draw=black] (0.8,0) -- (3.2,0) -- (3.2,2.4) -- (0.8,2.4) -- cycle;
		\filldraw[fill=orange, fill opacity=0.25, draw=black] (0,0) -- (3.2,0) -- (3.2,-1.8) -- (0,-1.8) -- cycle;
		\draw[line width=1pt] (0,0) -- (5,0) -- (3.2,2.4) -- cycle;
		\draw[red, line width=1pt] (3.2,2.4) -- (3.2,0);
		\draw[dashed] (5,0) arc (0:-90:1.8);
		\draw (1.7,1.5) node {$h_c^2$};
		\draw (1.6,-0.9) node {$p\, q$};
	\end{tikzpicture}
\end{minipage}

\section*{Mögliche Schwierigkeiten und Differenzierung}
\begin{itemize}[itemsep=0pt]
	\item \textbf{Konstruktion in 1.1:} Wer den Thaleskreis nicht sieht, erhält den Tipp \textit{Wo liegen alle Punkte, die mit $A$ und $B$ einen rechten Winkel bilden?}
	\item \textbf{Verwechslung der Sätze:} Höhensatz und Kathetensatz werden oft vermischt ($h^2=c\cdot p$). Die Übersichtstabelle und das Flächenbild helfen, ebenso die Frage \textit{Welche Seite wird quadriert?}
	\item \textbf{Höhe nicht auf die Hypotenuse:} Bei 3.1 f) zuerst mit Pythagoras $q$ bestimmen, erst dann den Höhensatz verwenden.
	\item \textbf{Schwächere Lernende:} In 1.1 nur die Tabellenzeilen 2--4 ausfüllen und direkt die Vermutung formulieren.
	\item \textbf{Stärkere Lernende:} Teil B von 1.1, Aufgaben 3.4--3.6; zusätzlich: Kehrsatz formulieren (\textit{Gilt $h_c^2=p\cdot q$, ist das Dreieck dann rechtwinklig?}).
\end{itemize}

\newpage
\section*{Lösungen}
\textbf{Aufgabe 1.1}
\begin{itemize}[itemsep=0pt]
	\item[a)--b)] $\overline{AB}=6.5$\,cm, $H_c$ mit $\overline{AH_c}=4.5$\,cm, Senkrechte in $H_c$ mit Thaleskreis schneiden. Messwert $h_c=3$\,cm.
	\item[c)] Zeile 1: $c=6.5$, $h_c^2=9=p\cdot q$. Zeile 2: $c=10$, $16=16$. Zeile 3: $c=12.5$, $25=25$. Zeile 4: $c=10$, $9=9$.
	\item[d)] Vermutung: $h_c^2=p\cdot q$.
	\item[e)--f)] $a^2=h_c^2+q^2$, $b^2=h_c^2+p^2$. Einsetzen: $2h_c^2+p^2+q^2=p^2+2pq+q^2$, also $h_c^2=p\cdot q$.
\end{itemize}

\textbf{Beweis Satz 2.1:} (I) $h_c^2=b^2-p^2$, (II) $b^2=c\cdot p$, also $h_c^2=c\cdot p-p^2=p\,(c-p)=p\cdot q$.

\textbf{Übersicht:} Pythagoras $a^2+b^2=c^2$ (Summe der Kathetenquadrate = Hypotenusenquadrat); Kathetensatz $a^2=c\cdot q$, $b^2=c\cdot p$ (Kathetenquadrat = Hypotenuse mal anliegender Abschnitt); Höhensatz $h_c^2=p\cdot q$ (Höhenquadrat = Produkt der Abschnitte).

\textbf{Beispiel 1:} $p=\frac{h_c^2}{q}=\frac{36}{4}=9$\,cm, $c=13$\,cm, $a=\sqrt{c\cdot q}=\sqrt{52}=2\sqrt{13}\approx 7.21$\,cm, $b=\sqrt{c\cdot p}=\sqrt{117}=3\sqrt{13}\approx 10.82$\,cm, $A=\frac{c\cdot h_c}{2}=39$\,cm$^2$. Kontrolle: $\frac{a\cdot b}{2}=\frac{6\cdot 13}{2}=39$\,cm$^2$.

\textbf{Aufgabe 3.1}
\begin{enumerate}[label=\alph*), itemsep=0pt]
	\item $h_c=\sqrt{16}=4$\,cm, $c=10$\,cm, $a=\sqrt{10\cdot 2}=2\sqrt5\approx 4.47$\,cm, $b=\sqrt{10\cdot 8}=4\sqrt5\approx 8.94$\,cm.
	\item $q=\frac{144}{16}=9$\,cm, $c=25$\,cm.
	\item $p=\frac{36}{3}=12$\,cm, $c=15$\,cm, $a=\sqrt{45}=3\sqrt5\approx 6.71$\,cm, $b=\sqrt{180}=6\sqrt5\approx 13.42$\,cm.
	\item $q=8$\,cm, $h_c=\sqrt{16}=4$\,cm, $b=\sqrt{10\cdot 2}=2\sqrt5\approx 4.47$\,cm.
	\item $h_c=\sqrt{25-9}=4$\,cm, $q=\frac{16}{3}\approx 5.33$\,cm, $c=\frac{25}{3}\approx 8.33$\,cm.
	\item $q=\sqrt{169-144}=5$\,cm, $p=\frac{144}{5}=28.8$\,cm, $c=33.8$\,cm.
\end{enumerate}

\textbf{Aufgabe 3.2:} $p+q=5q=15$, also $q=3$\,cm, $p=12$\,cm, $h_c=\sqrt{36}=6$\,cm, $A=\frac{15\cdot 6}{2}=45$\,cm$^2$.

\textbf{Aufgabe 3.3:} Nach Thales ist $\triangle ABC$ rechtwinklig bei $C$, $\overline{HC}$ ist die Höhe. Höhensatz: $\overline{HC}^2=6\cdot 2=12$\,cm$^2$ = Rechtecksfläche, also $\overline{HC}=\sqrt{12}\approx 3.46$\,cm.\\
e) Gleich: Thaleskreis und Senkrechte. Anders: Beim Höhensatz ist die Strecke $\overline{AB}$ die Summe der Rechteckseiten ($8$\,cm) und das Quadrat liegt über der Höhe. Beim Kathetensatz ist $\overline{AB}$ die längere Rechteckseite ($6$\,cm), die kürzere wird darauf abgetragen und das Quadrat liegt über einer Kathete.

\textbf{Aufgabe 3.4:} Es braucht $p\cdot q=10$ bzw.\ $p\cdot q=7$, z.B.\ a) $p=5$\,cm, $q=2$\,cm (oder $10$ und $1$), b) $p=7$\,cm, $q=1$\,cm (oder $3.5$ und $2$). Es gibt unendlich viele Möglichkeiten.

\textbf{Aufgabe 3.5:}
a) $p=1.5$\,m, $q=8.5$\,m: $h^2=12.75$, $h\approx 3.57$\,m.
b) Die Kante des Lastwagens liegt $5-1.25=3.75$\,m vom Rand: $h^2=3.75\cdot 6.25\approx 23.44$, $h\approx 4.84$\,m $>4.5$\,m, er passt.
c) $p\,(10-p)=16$ liefert $p=2$\,m (Probieren), also höchstens $10-2\cdot 2=6$\,m breit.

\textbf{Aufgabe 3.6:} a) Der Radius $\overline{MC}$ (mit Mittelpunkt $M$ der Strecke) ist das arithmetische Mittel $\frac{p+q}{2}$, die Höhe $\overline{HC}$ ist nach dem Höhensatz das geometrische Mittel $\sqrt{p\cdot q}$. b) Im rechtwinkligen Dreieck $\triangle MHC$ ist $\overline{MC}$ die Hypotenuse, also $\overline{HC}\leq\overline{MC}$. Gleichheit genau dann, wenn $H=M$, also $p=q$.

\section*{Quellen und Inspiration}
\begin{itemize}[itemsep=0pt]
	\item Geometrie 2 (Niederberger), Kap.\ \textit{Der Höhensatz}, S.\,83--84 (Herleitung, Aufgaben 1, 2, 5, 6, Vertiefung Nr.\,1).
	\item Geometrie 1 Aufgaben, Kap.\,5.2 \textit{Höhensatz und Kathetensatz}, Nr.\,40, 41, 43.
	\item Prüfungssammlung \texttt{SchoolTex/Exams} geprüft: Einzige thematisch passende Aufgabe (GZP25, Kathetenquadrate unterscheiden sich um $20$\,cm$^2$) wurde nicht verwendet.
\end{itemize}
\end{document}
